\documentclass[11pt]{article}

\usepackage{amsmath, amssymb}
\usepackage[margin=2.5cm]{geometry}
\usepackage{booktabs}
\usepackage{hyperref}
\usepackage{xcolor}

\newcommand{\param}[1]{\texttt{\color{blue!60!black}#1}}

\title{Synthetic Displacement Fields for DIC Image Generation \\
\large Equations and Parameters (matching \texttt{config.json})}
\author{}
\date{}

\begin{document}
\maketitle

\section*{Conventions}

All fields are defined on a normalized grid $X, Y \in [-1, 1]$, sampled on an
$h \times w$ pixel grid via \texttt{np.meshgrid(np.linspace(-1,1,w), np.linspace(-1,1,h))}.
Each mode contributes additively to the total field:
\[
u(X,Y) = \sum_{\text{active modes}} u_{\text{mode}}(X,Y), \qquad
v(X,Y) = \sum_{\text{active modes}} v_{\text{mode}}(X,Y)
\]
where $u = \texttt{disp\_x}$ and $v = \texttt{disp\_y}$. Parameter names in
\param{blue monospace} correspond exactly to keys in \texttt{config.json}.
Unless stated otherwise, $A \sim \mathcal{U}(0.1, \param{max\_disp})$ denotes an
amplitude redrawn independently for every sample and every mode instance.

\subsection*{Global parameters (sampled once per generated sample)}

\begin{align}
\nu &\sim \mathcal{U}(0, \param{nu\_max}) \\
\text{plane\_stress} &\sim \text{Bernoulli}(\param{plane\_stress\_probability})
\end{align}

$\nu$ (Poisson's ratio) is reused by the biaxial, necking/barreling transverse
coupling, and crack-tip modes. \texttt{plane\_stress} only affects the
crack-tip modes.

\subsection*{Final safety cap}

After all active modes are summed, if \param{max\_total\_disp} is not \texttt{null}:
\[
\|\mathbf{d}(X,Y)\| = \sqrt{u^2 + v^2}, \qquad
\mathbf{d}(X,Y) \leftarrow \mathbf{d}(X,Y) \cdot
\min\!\left(1,\ \frac{\param{max\_total\_disp}}{\|\mathbf{d}(X,Y)\|}\right)
\]

\newpage
\section{Rigid / Homogeneous Modes}

\subsection{Stretch (X, Y)}
\begin{align}
a &\sim \mathcal{U}(-1,1)\cdot \param{max\_disp} \\
u &\mathrel{+}= a\,X \quad \text{(\texttt{stretch\_x})} \\
v &\mathrel{+}= a\,Y \quad \text{(\texttt{stretch\_y})}
\end{align}

\subsection{Shear (X, Y)}
\begin{align}
a &\sim \mathcal{U}(-1,1)\cdot \param{max\_disp} \\
u &\mathrel{+}= a\,Y \quad \text{(\texttt{shear\_x})} \\
v &\mathrel{+}= a\,X \quad \text{(\texttt{shear\_y})}
\end{align}

\subsection{Translation}
\begin{align}
t_x, t_y &\sim \mathcal{U}(-\param{translation\_max}, \param{translation\_max}) \\
u &\mathrel{+}= t_x, \qquad v \mathrel{+}= t_y
\end{align}

\subsection{Rotation}
\begin{align}
\theta &\sim \mathcal{U}(-\param{rotation\_max}, \param{rotation\_max}) \\
u &\mathrel{+}= -\theta\,Y, \qquad v \mathrel{+}= \theta\,X
\end{align}

\subsection{Isotropic Dilation}
\begin{align}
\varepsilon &\sim \mathcal{U}(-1,1)\cdot \param{max\_disp} \\
u &\mathrel{+}= \varepsilon\,X, \qquad v \mathrel{+}= \varepsilon\,Y
\end{align}

\subsection{Poisson-Coupled Biaxial Stretch}
\begin{align}
\varepsilon_x &\sim \mathcal{U}(-1,1)\cdot \param{max\_disp} \\
u &\mathrel{+}= \varepsilon_x\,X \\
v &\mathrel{+}= -\nu\,\varepsilon_x\,Y
\end{align}

\section{Localized / Heterogeneous Modes}

\subsection{Localized Gaussian Displacement (\texttt{localized\_disp})}
Number of Gaussian sources $N \sim \mathcal{U}\{\param{n\_min}, \dots, \param{n\_max}\}$;
for each source $i = 1, \dots, N$:
\begin{align}
x_0, y_0 &\sim \mathcal{U}(-\param{pos\_range}, \param{pos\_range}) \\
\sigma &\sim \mathcal{U}(\param{sigma\_min}, \param{sigma\_max}) \\
G(X,Y) &= \exp\!\left(-\frac{(X-x_0)^2 + (Y-y_0)^2}{2\sigma^2}\right) \\
u &\mathrel{+}= A\,(X-x_0)\,G, \qquad v \mathrel{+}= A\,(Y-y_0)\,G
\end{align}

\subsection{Shear Band}
\begin{align}
\theta_{\text{band}} &\sim \mathcal{U}(0, \pi) \\
p &\sim \mathcal{U}(-\param{pos\_range}, \param{pos\_range}) \\
w &\sim \mathcal{U}(\param{width\_min}, \param{width\_max}) \\
B(X,Y) &= \exp\!\left(-\frac{(X\cos\theta_{\text{band}} + Y\sin\theta_{\text{band}} - p)^2}{2w^2}\right) \\
u &\mathrel{+}= A\,B\cos\theta_{\text{band}}, \qquad v \mathrel{+}= A\,B\sin\theta_{\text{band}}
\end{align}

\subsection{High Frequency}
\begin{align}
k &\sim \mathcal{U}(\param{k\_min}, \param{k\_max}) \\
\phi &\sim \mathcal{U}(0, 2\pi) \\
w(X,Y) &= A\sin\!\big(k(X+Y) + \phi\big) \\
u &\mathrel{+}= w, \qquad v \mathrel{+}= w
\end{align}

\subsection{Radial Displacement (image-centered)}
\begin{align}
A &\sim \mathcal{U}(-\param{max\_disp}, \param{max\_disp}) \\
R^2 &= X^2 + Y^2 \\
u &\mathrel{+}= A\,R^2, \qquad v \mathrel{+}= A\,R^2
\end{align}
\emph{Note:} centered on the image origin, grows as $R^2$ — distinct from
\texttt{cavity} below, which is a local $1/r$ point source at a random center.

\section{Damage / Fracture Modes}

\subsection{Necking (tension)}
\begin{align}
\sigma &\sim \mathcal{U}(\param{sigma\_min}, \param{sigma\_max}) \\
\text{band}(X) &= \exp\!\left(-\frac{X^2}{2\sigma^2}\right) \\
u &\mathrel{+}= A\,X\,\text{band}(X) \\
v &\mathrel{+}= -\nu\,A\,Y\,\text{band}(X)
\end{align}

\subsection{Barreling (compression)}
\begin{align}
\sigma &\sim \mathcal{U}(\param{sigma\_min}, \param{sigma\_max}) \\
\text{band}(Y) &= \exp\!\left(-\frac{Y^2}{2\sigma^2}\right) \\
u &\mathrel{+}= A\,X\,\text{band}(Y) \\
v &\mathrel{+}= -\nu\,A\,Y\,\text{band}(Y)
\end{align}

\subsection{Bending}
\begin{align}
\kappa &\sim \mathcal{U}(-\param{max\_disp}, \param{max\_disp}) \\
u &\mathrel{+}= -\kappa\,X\,Y \\
v &\mathrel{+}= \tfrac{1}{2}\kappa\,X^2
\end{align}

\subsection{Torsion}
\begin{align}
\alpha &\sim \mathcal{U}(-\param{max\_disp}, \param{max\_disp}) \\
R^2 &= X^2 + Y^2 \\
u &\mathrel{+}= -\alpha\,R^2\,Y, \qquad v \mathrel{+}= \alpha\,R^2\,X
\end{align}

\subsection{Cavity Expansion}
\begin{align}
x_0, y_0 &\sim \mathcal{U}(-\param{pos\_range}, \param{pos\_range}) \\
r &= \sqrt{(X-x_0)^2 + (Y-y_0)^2} \\
u &\mathrel{+}= A\,\frac{X - x_0}{r + \epsilon}, \qquad
v \mathrel{+}= A\,\frac{Y - y_0}{r + \epsilon} \qquad (\epsilon = 10^{-6})
\end{align}

\subsection{Inclusion}
\begin{align}
x_0, y_0 &\sim \mathcal{U}(-\param{pos\_range}, \param{pos\_range}) \\
\sigma &\sim \mathcal{U}(\param{sigma\_min}, \param{sigma\_max}) \\
G(X,Y) &= \exp\!\left(-\frac{(X-x_0)^2 + (Y-y_0)^2}{2\sigma^2}\right) \\
u &\mathrel{+}= A\,(X-x_0)\,G, \qquad v \mathrel{+}= A\,(Y-y_0)\,G
\end{align}

\subsection{Void Coalescence}
Cluster of $N \sim \mathcal{U}\{\param{n\_min},\dots,\param{n\_max}\}$ interacting cavities:
\begin{align}
c_x, c_y &\sim \mathcal{U}(-\param{cluster\_pos\_range}, \param{cluster\_pos\_range}) \\
s &\sim \mathcal{U}(\param{spread\_min}, \param{spread\_max})
\end{align}
For each void $i = 1, \dots, N$:
\begin{align}
x_{0,i} &= c_x + \delta_{x,i}, \quad \delta_{x,i} \sim \mathcal{U}(-s, s) \\
y_{0,i} &= c_y + \delta_{y,i}, \quad \delta_{y,i} \sim \mathcal{U}(-s, s) \\
A_i &\sim \mathcal{U}(0.1, \param{max\_disp}) / N \\
r_i &= \sqrt{(X-x_{0,i})^2 + (Y-y_{0,i})^2} \\
u &\mathrel{+}= A_i\,\frac{X - x_{0,i}}{r_i + \epsilon}, \qquad
v \mathrel{+}= A_i\,\frac{Y - y_{0,i}}{r_i + \epsilon}
\end{align}

\subsection{Crack Opening / Sliding (smoothed step)}
\begin{align}
x_0 &\sim \mathcal{U}(-\param{pos\_range}, \param{pos\_range}) &&
\text{(\texttt{crack\_open})} \\
y_0 &\sim \mathcal{U}(-\param{pos\_range}, \param{pos\_range}) &&
\text{(\texttt{crack\_slide})} \\
\delta &\sim \mathcal{U}(\param{delta\_min}, \param{delta\_max}) \\
v &\mathrel{+}= A\tanh\!\left(\frac{X - x_0}{\delta}\right) && \text{(opening)} \\
u &\mathrel{+}= A\tanh\!\left(\frac{Y - y_0}{\delta}\right) && \text{(sliding)}
\end{align}
A cheap qualitative discontinuity; \texttt{crack\_open} and \texttt{crack\_slide}
each have their own independent \param{pos\_range}/\param{delta\_min}/\param{delta\_max}
in \texttt{config.json}. For a physically accurate near-tip field, use the
LEFM modes below instead.

\subsection{Crack Tip Fields (LEFM, Williams series)}
Crack tip at random position, with crack line along the negative-$X$ axis
from the tip:
\begin{align}
x_0, y_0 &\sim \mathcal{U}(-\param{pos\_range}, \param{pos\_range}) \\
r &= \sqrt{(X-x_0)^2 + (Y-y_0)^2} + \epsilon \\
\theta &= \operatorname{atan2}(Y - y_0,\ X - x_0) \\
\kappa &=
\begin{cases}
\dfrac{3-\nu}{1+\nu}, & \text{plane\_stress} = \text{True} \\[1.5ex]
3 - 4\nu, & \text{plane\_stress} = \text{False}
\end{cases}
\end{align}

\textbf{Mode I (opening), \texttt{crack\_tip\_KI}:}
\begin{align}
K_I &\sim \mathcal{U}(0.1, \param{max\_disp}) \\
u &\mathrel{+}= K_I \sqrt{\frac{r}{2\pi}} \cos\!\frac{\theta}{2}
\left[\kappa - 1 + 2\sin^2\!\frac{\theta}{2}\right] \\
v &\mathrel{+}= K_I \sqrt{\frac{r}{2\pi}} \sin\!\frac{\theta}{2}
\left[\kappa + 1 - 2\cos^2\!\frac{\theta}{2}\right]
\end{align}

\textbf{Mode II (in-plane shear), \texttt{crack\_tip\_KII}:}
\begin{align}
K_{II} &\sim \mathcal{U}(0.1, \param{max\_disp}) \\
u &\mathrel{+}= K_{II} \sqrt{\frac{r}{2\pi}} \sin\!\frac{\theta}{2}
\left[\kappa + 1 + 2\cos^2\!\frac{\theta}{2}\right] \\
v &\mathrel{+}= -K_{II} \sqrt{\frac{r}{2\pi}} \cos\!\frac{\theta}{2}
\left[\kappa - 1 - 2\sin^2\!\frac{\theta}{2}\right]
\end{align}
Both modes share the same random tip position $(x_0, y_0)$ when both are
active in the same sample. $K_I$, $K_{II}$ are synthetic pixel-scale
amplitudes (not physical MPa$\sqrt{\text{m}}$ stress intensity factors).

\subsection{Buckling}
Number of harmonics $N_h \sim \mathcal{U}\{\param{n\_harmonics\_min}, \dots, \param{n\_harmonics\_max}\}$;
for each harmonic $i = 0, \dots, N_h-1$:
\begin{align}
A_i &\sim \mathcal{U}(0.1, \param{max\_disp}) / (i+1) \\
k_i &\sim \mathcal{U}(\param{k\_min}, \param{k\_max}) \cdot (i+1) \\
\phi_i &\sim \mathcal{U}(0, 2\pi) \\
v &\mathrel{+}= \sum_{i=0}^{N_h-1} A_i \sin\!\big(\pi k_i X + \phi_i\big)
\end{align}

\subsection{Contact / Indentation (Hertzian-inspired)}
\begin{align}
x_0 &\sim \mathcal{U}(-\param{pos\_range}, \param{pos\_range}) \\
a &\sim \mathcal{U}(\param{a\_min}, \param{a\_max}) \\
P(X) &= \exp\!\left(-\frac{(X-x_0)^2}{2a^2}\right) \\
v &\mathrel{+}= -A\,P(X) \\
u &\mathrel{+}= 0.5\,A\,\frac{X - x_0}{a}\,P(X)
\end{align}

\subsection{Delamination / Blister}
\begin{align}
x_0, y_0 &\sim \mathcal{U}(-\param{pos\_range}, \param{pos\_range}) \\
\sigma &\sim \mathcal{U}(\param{sigma\_min}, \param{sigma\_max}) \\
r &= \sqrt{(X-x_0)^2 + (Y-y_0)^2} \\
\text{ring}(r) &= \frac{r}{\sigma}\exp\!\left(-\frac{r^2}{2\sigma^2}\right) \\
u &\mathrel{+}= A\,\frac{X-x_0}{r+\epsilon}\,\text{ring}(r), \qquad
v \mathrel{+}= A\,\frac{Y-y_0}{r+\epsilon}\,\text{ring}(r)
\end{align}
Ring-shaped: zero at the blister center, peaks at $r \approx \sigma$ (the
delamination front), decays beyond — distinct from \texttt{inclusion}
(peaks at center) and \texttt{cavity} (monotonic $1/r$ decay).

\section*{Parameter Reference Table}

\begin{center}
\begin{tabular}{@{}ll@{}}
\toprule
\textbf{JSON key} & \textbf{Used by} \\
\midrule
\param{max\_disp} & all modes (amplitude budget) \\
\param{max\_total\_disp} & final magnitude cap (all modes) \\
\param{rotation\_max} & \texttt{rotation} \\
\param{translation\_max} & \texttt{translation} \\
\param{nu\_max} & \texttt{poisson\_biaxial}, \texttt{necking}, \texttt{barreling}, crack-tip modes \\
\param{plane\_stress\_probability} & \texttt{crack\_tip\_KI}, \texttt{crack\_tip\_KII} \\
\param{min\_modes}, \param{max\_modes} & mode selection (how many combined per sample) \\
\bottomrule
\end{tabular}
\end{center}

Per-mode parameter blocks (\param{localized\_disp}, \param{shear\_band},
\param{high\_freq}, \param{necking}, \param{barreling}, \param{cavity},
\param{inclusion}, \param{void\_coalescence}, \param{crack\_open},
\param{crack\_slide}, \param{crack\_tip}, \param{buckling},
\param{contact\_indentation}, \param{delamination\_blister})
each contain only the ranges relevant to that mode, exactly as listed
in the equations above.

\end{document}
